% Complete chapter-1 typesetting source; compact textbook presentation.
% Authoritative Markdown: ../../chapters/ch01/answer-key.md
% Stable task/source IDs are comments. All learner prose, tasks, and notes are retained.
% Shared layout for the two complete chapter-1 PDFs.
% This is a small hand-maintained TeX layout, not a conversion pipeline.
\RequirePackage{fix-cm}
\documentclass[11pt,twoside]{article}
\usepackage{fontspec}
\setmainfont{Cambria}[
  Renderer=HarfBuzz,
  BoldFont=cambriab.ttf,
  ItalicFont=cambriai.ttf,
  BoldItalicFont=cambriaz.ttf
]
\setmonofont{lmmono10-regular.otf}[BoldFont=lmmonolt10-bold.otf]
\usepackage[paperwidth=165mm,paperheight=241mm,inner=21mm,outer=19mm,
  top=22mm,bottom=20mm,headheight=14pt,headsep=7mm,footskip=10mm]{geometry}
\usepackage{xcolor,array,booktabs,longtable,tabularx}
\usepackage{enumitem,fancyhdr}
\usepackage[unicode,hidelinks]{hyperref}
\definecolor{ink}{HTML}{000000}
\definecolor{muted}{HTML}{4C5558}
\setlength{\parindent}{0pt}
\setlength{\parskip}{3.5pt}
\AtBeginDocument{\fontsize{11}{14.3}\selectfont}
\setlength{\tabcolsep}{4pt}
\renewcommand{\arraystretch}{1.10}
\setlist[itemize]{leftmargin=1.3em,label=\textbullet,itemsep=2pt,topsep=3pt,parsep=0pt}
\setlist[enumerate]{leftmargin=1.8em,itemsep=3pt,topsep=4pt,parsep=0pt}
\widowpenalty=10000
\clubpenalty=10000
\emergencystretch=1em
\raggedbottom
\pagestyle{fancy}
\fancyhf{}
\fancyhead[LE,RO]{\small\itshape Naming the world}
\fancyhead[LO,RE]{\small Cook Expanded}
\fancyfoot[LE,RO]{\small\thepage}
\renewcommand{\headrulewidth}{0.3pt}
\newcommand{\booktitle}[2]{
  {\fontsize{9}{11}\selectfont\color{muted}#1\par}
  \vspace{3pt}
  {\fontsize{20}{24}\selectfont\color{ink}#2\par}
  \vspace{6pt}
}
\newcommand{\source}[1]{\par{\footnotesize\color{muted}#1\par}}
\newcommand{\readingtext}[1]{
  \par\addvspace{5pt}
  \begingroup\leftskip=1.5em
    \fontsize{12}{16}\selectfont #1\par
  \endgroup\addvspace{5pt}
}
\makeatletter
\renewcommand{\section}{\@startsection{section}{1}{\z@}
  {-2.4ex \@plus -0.6ex \@minus -0.2ex}{0.8ex \@plus 0.2ex}
  {\normalfont\fontsize{12}{15}\selectfont\bfseries\color{ink}}}
\renewcommand{\subsection}{\@startsection{subsection}{2}{\z@}
  {-1.8ex \@plus -0.5ex \@minus -0.2ex}{0.5ex \@plus 0.2ex}
  {\normalfont\fontsize{11}{14}\selectfont\bfseries\color{ink}}}
\makeatother

\hypersetup{pdftitle={Chapter 1: Answers and commentary},pdfauthor={Cook Expanded}}
\fancyfoot[LO,RE]{\footnotesize\color{muted}Answers and commentary}
\hyphenation{heofon heofone}
\begin{document}
\booktitle{CHAPTER 1 · ANSWER KEY}{Answers and commentary}


% Chapter: ch01. Compact presentation revision awaiting review. Stable item IDs match coverage.json.


\pdfbookmark[1]{Preparation (§5)}{section1}
\section*{Preparation (§5)}


\begin{enumerate}
% item: ch01-prep-01
\item[1.] \textbf{sēo eorðe}, “the earth”: feminine nominative singular. \textbf{sēo} alone fills the blank.
% item: ch01-prep-02
\item[2.] \textbf{ðæt ǣfen}, “the evening”: neuter. Masculine \textbf{se ǣfen} is grammatical but differs from the selected gender required here.
% item: ch01-prep-03
\item[3.] \textbf{ðā ðīestru}, “the darkness”: plural; \textbf{ðā} is nominative plural. The blank requires \textbf{ðā}, with the number stated.
\end{enumerate}


\pdfbookmark[1]{Composition (§6)}{section2}
\section*{Composition (§6)}


The target is two definite nominative noun phrases joined by \textbf{and}. Each article agrees with its own noun: \textbf{se} masculine singular, \textbf{sēo} feminine singular, \textbf{ðæt} neuter singular, \textbf{ðā} plural. This agreement explains the forms in every answer below; coordination alone does not require plural articles.


\begingroup\small
\begin{longtable}{@{}>{\raggedright\arraybackslash}p{8mm} >{\raggedright\arraybackslash}p{63mm} >{\raggedright\arraybackslash}p{48mm}@{}}
\toprule
\textbf{No.} & \textbf{Old English} & \textbf{Meaning} \\
\midrule
\endfirsthead
\multicolumn{3}{@{}l}{\small\itshape Continued} \\
\toprule
\textbf{No.} & \textbf{Old English} & \textbf{Meaning} \\
\midrule
\endhead
\midrule
\multicolumn{3}{r@{}}{\footnotesize Continued on the next page} \\
\endfoot
\bottomrule
\endlastfoot
% item: EX-I-01-01
1 & \textbf{ðæt ǣfen and se morgen} & The evening and the morning. \\
% item: EX-I-01-02
2 & \textbf{se dæg and sēo niht} & The day and the night. \\
% item: EX-I-01-03
3 & \textbf{sēo eorðe and se heofon} & The earth and the heaven. \\
% item: EX-I-01-04
4 & \textbf{se fisc and se fugol} & The fish and the bird. \\
% item: EX-I-01-05
5 & \textbf{ðæt gærs and ðæt sǣd} & The grass and the seed. \\
% item: EX-I-01-06
6 & \textbf{ðæt trēo and sēo wyrt} & The tree and the herb. \\
% item: EX-I-01-07
7 & \textbf{sēo drȳgnes and ðæt wæter} & The dry land and the water. \\
% item: EX-I-01-08
8 & \textbf{sēo andlīcnes and sēo gelīcnes} & The image and the likeness. \\
% item: EX-I-01-09
9 & \textbf{ðæt lēoht and ðā ðīestru} & The light and the darkness. \\
% item: EX-I-01-10
10 & \textbf{ðæt tācen and sēo tīd} & The sign and the season. \\
% item: EX-I-01-11
11 & \textbf{se gāst and ðæt līf} & The spirit and the life. \\
% item: EX-I-01-12
12 & \textbf{sēo gesceaft and se męte} & The creature and the food. \\
% item: EX-I-01-13
13 & \textbf{ðæt gēar and se dæg} & The year and the day. \\
% item: EX-I-01-14
14 & \textbf{se mann and ðæt dēor} & The man and the animal. \\
% item: EX-I-01-15
15 & \textbf{ðæt trēo and se wæstm} & The tree and the fruit. \\
% item: EX-I-01-16
16 & \textbf{sēo fæstnes and sēo nēowolnes} & The firmament and the abyss. \\
\end{longtable}
\endgroup


\clearpage
\textbf{Variants.} In 1, masculine \textbf{se ǣfen}; in 3, weak feminine \textbf{sēo heofone}; in 6 and 15, \textbf{ðæt trēow}; in 8, \textbf{sēo anlīcnes}; in 12, neuter \textbf{ðæt gesceaft}. These are grammatical alternatives with the required meanings, though they differ from the selected vocabulary forms or genders. \textbf{heofone}, not masculine \textbf{heofon}, takes the feminine article in that alternative.


\textbf{Meaning and form.} In 7, \textbf{drȳgnes} is one noun, “dry land.” In 9, only \textbf{ðīestru} is plural. \textbf{męte} in 12 means food generally; \textbf{dēor} in 14 means animal, not specifically deer. In 15, \textbf{wæstm} has short \textbf{æ}.


\pdfbookmark[1]{Reading (§7)}{section3}
\section*{Reading (§7)}


\begin{enumerate}
% item: ch01-reading-01
\item[1.] The earth; it was empty or void. Either adjective expresses emptiness.
% item: ch01-reading-02
\item[2.] \textbf{Sēo} (article) + \textbf{eorðe} (noun): feminine nominative singular, subject of \textbf{wæs}. \textbf{and} joins the predicate adjectives \textbf{īdel} and \textbf{ǣmtig}.
% item: ch01-reading-03
\item[3.] “The earth, for its part, was empty and void.” Also: “Now the earth was empty and vacant.” \textbf{wæs} requires past tense; \textbf{sōðlice} is connective here. The two adjectives overlap in meaning: no unique pair of English synonyms is required.
\end{enumerate}


“Indeed” is also possible for \textbf{sōðlice}, though the connective is preferred in context. A single expression such as “completely empty” conveys the overlapping adjectives; a close rendering should represent both. A final period is acceptable for the extracted clause. Present-tense “is” misreads \textbf{wæs}; “without form” is a freer contextual rendering, not a required gloss.


\pdfbookmark[1]{Review (§8)}{section4}
\section*{Review (§8)}


\begin{enumerate}
% item: ch01-review-01
\item[1.] \textbf{ðæt sǣd and sēo wyrt}, “the seed and the herb”: neuter and feminine nominative singular. Each article agrees with its own noun.
% item: ch01-review-02
\item[2.] \textbf{se heofon and ðæt wæter}, “the heaven and the water”: masculine and neuter nominative singular. \textbf{sēo heofone and ðæt wæter} is grammatical with the feminine variant, though \textbf{heofon} is the recall target.
% item: ch01-review-03
\item[3.] \textbf{ðā} agrees with plural \textbf{ðīestru}; \textbf{ðæt} with neuter singular \textbf{lēoht}. The plural article is the same for all genders. One plural noun form supplies no general rule for the other nouns.
\end{enumerate}


\clearpage
\pdfbookmark[1]{Assessment notes}{section5}
\section*{Assessment notes}


Distinguish grammatical validity, correct meaning, and compliance with the specified vocabulary and construction. A verified synonym or another case can be valid Old English while missing the exercise's target. Equivalent explanations and grammatical abbreviations are sufficient for the analytical questions.


The common agreement errors are choosing gender from English meaning, using one article form despite different noun genders, or treating the English singular “darkness” as evidence that \textbf{ðīestru} is singular. Use the gender and number of each Old English noun.


Capitalization and final punctuation do not determine grammatical correctness. Thorn/eth alternatives and emphatic \textbf{sē} beside ordinary \textbf{se} are acceptable. Missing vowel marks are spelling errors to correct separately from syntax.


\begingroup\small
\pdfbookmark[1]{Notes on sources and variants}{section6}
\section*{Notes on sources and variants}


\textbf{Ex:} A. S. Cook, \textit{Exercises in Old English} (1895). \textbf{FB:} \textit{A First Book in Old English}, 2nd ed. (1895). References use printed pages. All answers and translations are newly prepared; full item evidence is in the editorial record.


\begin{enumerate}
\item[1.] \textbf{Preparation, composition, review:} Ex I.1, p.1, and vocabulary; FB §§43–44, 46–47, 51–53; articles §84, p.50; thorn/eth §§2–3, p.5. Plural \textbf{ðīestru}: Ex p.44; FB p.123, note 7.
\item[2.] \textbf{Variants:} \textbf{ǣfen}, FB §47.7, p.33, and \href{https://bosworthtoller.com/37608}{Bosworth–Toller supplement}; \textbf{heofon/\allowbreak heofone}, FB §53.3, p.37; \textbf{trēo/trēow}, §47.3, p.33; \textbf{an(d)līcnes}, Ex p.51; feminine/neuter \textbf{gesceaft}, \href{https://bosworthtoller.com/15847}{Bosworth–Toller}.
\item[3.] \textbf{Correction C01:} Short-vowel \textbf{wæstm} replaces Cook's \textbf{wǣstm}: \href{https://glossary.oldenglishaerobics.net/oeaginfo.php?objid=xwaestm\&objtype=gentry}{Baker's glossary}; \href{https://doi.org/10.17811/selim.29.2024.109-25}{Stenroos (2024), p.118, Table 1}.
\item[4.] \textbf{Reading:} FB I, p.123, second sentence's opening clause and notes 4–6; note 5 gives \textit{autem}. \textbf{wæs}, §138, pp.78–79; subject/predicate, §§149–150, p.88; adjective forms, §57.3, p.39; vocabulary pp.274, 301, 313.
\end{enumerate}



\endgroup
\end{document}
